\documentclass[a4paper,10pt]{article}
\usepackage[utf8]{inputenc}
\usepackage{amsmath}
\usepackage{amsfonts}
\usepackage{mathtools}
\title{Range Proofs Using Linear Constraints}
\author{Armando Faz-Hernandez}
\date{August 17, 2026}
\begin{document}
\maketitle
Let $x\in \mathbb{Z}$ such that $x\geq 0$,  $x$ can be written in base-$B$, with $B\geq2$, using $n=\lfloor \log_B(x)\rfloor+1$ digits. So $(d_{n-1},\dots,d_{0})_B$ is the representation of $x$ in base-$B$ if it holds
$$ x= \sum_{i=0}^{n-1} d_i B^i, \text{ where } d_i\in\{0,\dots, B-1\}\,.$$

\paragraph{Valid Digit}
We want to build linear constraints to check that $d\in\{0,\dots,k\}$ for $k>0$.
Let $H_1,H_2$ be generators of a group $\mathbb{G}$.
First, we commit to the digit value $C_0 = \text{Comm}(d\,;r)$, with $r$ be a random scalar. Using a Pedersen commitment, we have
$$ C_0 = dH_1 + rH_2\,.$$
Then, we define the following values:
$$ C_{i} \coloneq (d-i)C_{i-1} - r_i H_2, \text{ where } r_i = r\prod_{j=1}^{i} (d-j), \text{ for } 1 \leq i \leq k\,.$$
Every equation adds a factor $(d-i)$ that builds a polynomial whose roots are $i$ for all $i\in\{0,\dots,k\}$.

We redistribute the equations, so the witness for $d$ is $w\coloneq(d,r,r_1,\dots,r_k)$, and the linear relations are:
$$
\begin{aligned}
C_0      &= dH_1 + sH_2\\
iC_{i-1} &= dC_{i-1} - r_i H_2, \text{ for } 1 \leq i \leq k\,.
\end{aligned}
$$

\paragraph{Positive Integer in Range}
We want to check that $x$ is in the range $0\leq x \leq N$ for some public $N$.
To do so, we write $x$ in base-$B$, for some $B$, and obtain $(d_{n-1},\dots,d_0)$, where $n=\lfloor \log_B(N)\rfloor+1$.
Then, to prove $x$ is in the range, we can apply the ValidDigit procedure to the digits distinguishing two cases.
\begin{enumerate}
 \item If $B^n-1=N$, apply ValidDigit to all digits setting $k=B-1$.
 \item Otherwise, apply ValidDigit to the digits $d_{i}$, such that $i<n-1$, setting $k=B-1$; and for the digit $d_{n-1}$, apply ValidDigit with $k=\lfloor N/B^{n-1} \rfloor$.
\end{enumerate}

\paragraph{Negative Integers}
We want to check that $x$ is in the range $-N \leq x \leq N$ for some public value $N$.
Set $s=1$ if $x<0$, otherwise, set $s=0$.
Apply the PositiveInteger procedure from above to $|x|$, and apply ValidDigit to $s$ setting $k=1$.

\end{document}
